\section{Integrationsmethoden}

\subsection{Substitution}
\textbf{Verwende bei:}
\[\int f(g(x))\cdot g'(x)\,dx\]
Wenn ein innerer Ausdruck \(g(x)\) und dessen Ableitung \(g'(x)\) als Faktor vorkommen.

\textbf{Vorgehen:}
\begin{enumerate}
    \item Substitutionsgleichung für
          \[
              x: u = g(x)
          \]
    \item Substitutionsgleichung für
          \[
              dx: \frac{du}{dx} = g'(x) \Rightarrow dx = \frac{du}{g'(x)}
          \]
    \item Integralsubstitution:
          \begin{align*}
              \int f\!\bigl(g(x)\bigr)\cdot g'(x)\,dx     & = \int f(u)\,du, \quad \text{unbestimmt}             \\
              \int_a^b f\!\bigl(g(x)\bigr)\cdot g'(x)\,dx & = \int_{g(a)}^{g(b)} f(u)\,du, \quad \text{bestimmt}
          \end{align*}
    \item  Rücksubstitution (nur für unbestimmtes Integral):
          \[
              \int f(u)\,du = F(u) + C = F\!\bigl(g(x)\bigr) + C
          \]

\end{enumerate}

\textbf{Spezialfall:}\\
Für eine Funktion \(f\) und Konstanten \(a, b\) gilt:
\[
    \int f(ax+b)\,dx = \frac{1}{a} F(ax+b) + C
\]

\subsection{Partielle Integration}
\textbf{Verwende bei:}
\[\int u(x)\cdot v'(x)\,dx\]
Bei Produkten zweier Funktionen, insbesondere wenn ein Faktor sich durch Ableiten vereinfacht (z.\,B.\ \(x\cdot e^x\), \(x\cdot\sin x\), \(\ln x\)).

\textbf{Formel:}
\begin{align*}
    \int u(x)\cdot v'(x)\,dx &= u(x)\cdot v(x) - \int u'(x)\cdot v(x)\,dx \\
    \int_a^b u(x)\cdot v'(x)\,dx &= \bigl[u(x)\cdot v(x)\bigr]_a^b - \int_a^b u'(x)\cdot v(x)\,dx
\end{align*}

\textbf{Faustregel (LIATE):} Wähle \(u\) in dieser Reihenfolge:
\begin{enumerate}
    \item[\textbf{L}] Logarithmen (\(\ln x\))
    \item[\textbf{I}] Inverse Trig.\ (\(\arctan x\), \(\arcsin x\))
    \item[\textbf{A}] Algebraisch (\(x^n\))
    \item[\textbf{T}] Trigonometrisch (\(\sin x\), \(\cos x\))
    \item[\textbf{E}] Exponentiell (\(e^x\))
\end{enumerate}

\subsection{Partialbruchzerlegung}
\textbf{Verwende bei:} Integralen rationaler Funktionen 
\[
\int f(x)\,dx = \int \frac{P(x)}{Q(x)}\,dx
\]
Falls Grad \(P \geq Q\), zuerst Polynomdivision durchführen.

\textbf{Vorgehen:}
\begin{enumerate}
    \item Reelle Nullstellen von \(Q(x)\) bestimmen und \(Q(x)\) in Linearfaktoren zerlegen.
    \item Partialbrüche aufstellen je nach Faktorgrad:
          \begin{itemize}
              \item Einfache reelle Nullstelle \((x - a)\): \(\dfrac{A}{x-a}\)
              \item \(k\)-fache reelle Nullstelle \((x-a)^k\): \(\dfrac{A_1}{x-a} + \dfrac{A_2}{(x-a)^2} + \cdots + \dfrac{A_k}{(x-a)^k}\)
          \end{itemize}
    \item \(f(x)\) mit Summe der Partialbrüche gleichsetzen und auf gemeinsamen Nenner bringen.
    \item Konstanten \(A_i\) durch einsetzen bestimmen.
    \item Einzelne Brüche integrieren:
          \begin{align*}
              \int \frac{1}{x-x_0}\,dx     & = \ln|x-x_0|+C                   \\
              \int \frac{1}{(x-x_0)^r}\,dx & = \frac{1}{(1-r)(x-x_0)^{r-1}}+C \\
          \end{align*}
\end{enumerate}
